\section{The ETH lower bound}
\label{sec:complexity}

The parity reduction is unconditional. To obtain its complexity
consequences we first control domain sizes, then account for the
sparsification and graph-construction costs before choosing the dimension.
We use the width-three case of the sparsification lemma
\citep[Corollary~1, p.~523]{IPZ2001}.
Here clauses have at most three literals, and \(|F|\) denotes the explicit
encoding length of a formula \(F\).

\begin{lemma}[Sparsification]
\label{lem:sparsification}
For every \(\epsilon>0\) there are a constant \(C_\epsilon\ge1\) and a
deterministic algorithm that expresses any \(3\)-CNF \(F\) on \(N\ge1\)
variables as a disjunction of at most \(2^{\epsilon N}\) \(3\)-CNFs on
the same variables, each with at most \(C_\epsilon N\) clauses.
The running time is \(2^{\epsilon N}\operatorname{poly}(|F|)\).
For positive rational \(\epsilon\), the procedure and a valid density
bound can be chosen effectively.
\end{lemma}

The lemma is the imported sparsification theorem, not a consequence of
the parity construction. In its constructive set-system proof, the
universe consists of the \(2N\) literals: choosing rate \(\epsilon/2\)
there gives the displayed variable-based rate \(\epsilon\). The integer
threshold choice and fixed branching procedure provide the effective
rational-rate selection used below~\citep[pp.~520--523]{IPZ2001}.

The stated input-length bound allows an initial polynomial-time
normalization: rename variables and remove tautologies and repeated
clauses. There are only \(O(N^3)\) distinct clauses of width at most three,
so the normalized encoding has length polynomial in \(N\).

\begin{lemma}[Grouping sparse clauses]
\label{lem:sat-choice}
A \(3\)-CNF with \(m\) clauses yields a simultaneous-choice instance with
\(t\) domains, each domain and pair-label set having size at most
\[
A=2^{3\lceil m/t\rceil}.
\]
A simultaneous choice exists if and only if the formula is satisfiable.
\end{lemma}
\begin{proof}
Partition the clauses into \(t\) groups of size at most \(\lceil m/t\rceil\).
Let \(V_i\) be the variables in group \(i\), and let \(D_i\) contain exactly
the assignments on \(V_i\) satisfying that group. Set
\[
L_{ij}=\{0,1\}^{V_i\cap V_j},
\qquad
\lambda_{ij}^i(d)=d|_{V_i\cap V_j},
\qquad
\lambda_{ij}^j(d)=d|_{V_i\cap V_j}.
\]
Each scope has size at most \(3\lceil m/t\rceil\), proving the bounds.
Pairwise agreement assigns every shared variable one value, so the
choices join to an assignment satisfying all clauses; unused variables
are free. Conversely, restrict a satisfying assignment to each scope.
An empty group has its unique empty assignment, an unsatisfiable group
has empty domain, and an empty overlap has one label, the empty
assignment. Thus these cases require no exception.
\end{proof}

Apply Theorem~\ref{thm:choice-reduction} with \(t=k+1\) for the
joint convention, or \(t=k\) for separate coordinate multisets.
For a sparse formula with \(m\le CN\), recall that
\[
M_t=\sum_{a\in V(T)}2^{\deg_T(a)-1}=2^{O(k^2)}.
\]
There are at most \(A\) vertices in each base fiber, and precisely
\(2^{\deg_T(a)-1}\) tags per vertex of type \(a\). The two output graphs
therefore have a common order satisfying
\begin{equation}
n\le M_t\,2^{3\lceil m/t\rceil}
   \le b_k\,2^{3CN/t},\qquad b_k=8M_t.
\label{eq:output-size}
\end{equation}
These are explicit graphs. Enumerate the scope and overlap assignments,
retain the satisfying domain assignments, and enumerate the parity
tags. A scan of all vertex pairs then writes both adjacency matrices:
each entry is determined by restrictions of assignments and a comparison
of tag bits. The vertex records have length \(O(N+k^2)\), so a direct
multitape implementation takes at most
\begin{equation}
b'_k\,(N+|F|)^{d}\,2^{B_0 CN/t},
\qquad B_0=6,
\label{eq:construction-time}
\end{equation}
for an absolute \(d\) and a constant \(b'_k\) depending only on \(k\).
Indeed, squaring the assignment bound gives
\(2^{6\lceil m/t\rceil}\le64\,2^{6CN/t}\).
A fixed-power tape simulation only changes the absolute constant \(B_0\).

\begin{proof}[Proof of Theorem~\ref{thm:eth}]
Use positive-rate ETH \citep{IP2001}: for some \(\delta>0\), no
deterministic \(3\)-SAT algorithm runs in
\(2^{\delta N}\operatorname{poly}(|F|)\) time.
Choose \(\epsilon<\delta/4\), and fix \(C=C_\epsilon\) from
Lemma~\ref{lem:sparsification}.
For the Turing-machine model set \(q=1\). More generally, suppose the
chosen machine model has a simulation taking at most
\(A_k(T+\ell+1)^q\) steps, where \(q\ge1\) is fixed by the model,
independently of the simulated program and \(k\).
Here \(T\) is the program's running time and \(\ell\) its input length;
\(A_k\) may depend on the chosen fixed program and \(k\). Choose
\[
c=\frac{\delta}{24Cq},
\qquad
K\ge4,\qquad
K>\frac{4C\max\{B_0,6q\}}{\delta}.
\]

Fix any \(k\ge K\) and suppose a solver takes \(O(n^{ck})\) time.
Hardcode this solver and \(k\) into a \(3\)-SAT algorithm.
Sparsify the input, form the graph pair for every sparse disjunct, and
accept exactly when at least one pair is inequivalent.
Lemmas~\ref{lem:sparsification} and~\ref{lem:sat-choice}, together with
Theorem~\ref{thm:choice-reduction}, prove correctness.

The matrix input has length \(\ell=O(n^2)\), so the simulated solver
takes \(O_k(n^{qck}+n^{2q}+1)\) time. Both terms matter, even when
\(ck<2\). By \eqref{eq:output-size} and \eqref{eq:construction-time},
the total time is bounded, up to polynomial input factors and
fixed-\(k\) constants, by
\[
2^{\epsilon N}
\left(
  2^{B_0 CN/t}
 +2^{6qCN/t}
 +2^{3CqckN/t}
\right).
\]
Since \(t\ge k\), the first two inner rates are below \(\delta/4\),
and the third is at most \(\delta/8\).
The resulting rate is below \(\delta/2\), contradicting ETH.
The solver's description, its running-time constant, \(A_k\), \(b_k\),
and the eventual threshold in \(N\) may all depend on this fixed \(k\);
\(c\) and \(K\) do not. Thus the contradiction applies to every fixed
\(k\ge K\), without a uniform way of selecting the alleged solvers.
\end{proof}

\begin{remark}[Small dimensions on matrix inputs]
On a comparison of two edgeless \(n\)-vertex graphs, a deterministic algorithm must
inspect every potential edge: changing an unseen edge in one graph
destroys equivalence, by degrees for ordinary \(1\)-WL and by atomic
tuple types for \(k\ge2\). This gives an \(\Omega(n^2)\) bit-query bound.
Decreasing \(c\) so that \(ck<2\) for the finitely many
\(2\le k<K\) therefore includes those dimensions for both tuple
conventions. Ordinary neighbor color refinement at dimension one is
covered separately by the degree argument. The literal one-coordinate
replacement rules above, whose atomic types see no edges, are not this
ordinary \(1\)-WL convention.
\end{remark}

\subsection{The uniform subexponential formulation}
\label{sec:uniformity}

The conclusion under an assumption that only excludes a single uniform
subexponential \(3\)-SAT algorithm is explicitly uniform.
Here a uniform WL algorithm takes \(k\) and the two graphs as input.

\begin{proposition}[Uniform lower bound]
\label{prop:uniform}
If no deterministic \(3\)-SAT algorithm has running time
\(2^{o(N)}\operatorname{poly}(|F|)\), there is no uniform deterministic algorithm for
\(k\)-WL equivalence with running time \(f(k)n^{g(k)}\), where
\(g(k)/k\to0\).
\end{proposition}
\begin{proof}
The constructive sparsification proof chooses integer thresholds by
a computable search from any positive rational rate and then runs a
fixed branching procedure \citep[pp.~520--523]{IPZ2001}.
We may therefore fix an effectively indexed family \(S_j\) with rate
\(1/j\) and sparsity constant \(C_j\), for \(j\ge1\).

Suppose the uniform WL algorithm exists. We obtain one subexponential
solver by racing all fixed parameter choices, so no effective rate of
convergence for \(g(k)/k\) is needed. One fixed Turing machine
\(U(j,k,F)\) runs \(S_j\), constructs the graph pairs, and calls that
algorithm. For every fixed \(j\ge1\) and \(k\ge4\), its full running time
has the form
\[
T_{j,k}(F)\le A_{j,k}N^{d_j}2^{\rho_{j,k}N},
\qquad
\rho_{j,k}=\frac1j+
 \max\left\{\frac{B_0C_j}{t},\frac{3C_jg(k)}{t}\right\},
\]
after normalization. Parameter initialization and graph output are
included in \(T_{j,k}\).

Run doubling rounds \(R=1,2,4,\ldots\). In each round simulate \(U(j,k,F)\)
from scratch for at most \(R\) steps for every \(1\le j\le R\) and
\(4\le k\le R\), stopping at the first completed decision.
All these computations are correct. If \(L\) is the normalized input
length, a fixed tape simulator costs at most
\(O((R+L+1)^3)\) per trial, including input copying and resetting the
at most \(O(L+R)\) touched cells. There are at most \(R^2\) trials.
Consequently the single resulting machine has running time
\[
D(F)\le C_U\bigl(T_{j,k}(F)+L+j+k+4\bigr)^5
\]
for every fixed pair, with an absolute power and a constant \(C_U\)
independent of that pair.

Given \(\eta>0\), first fix \(j\) with \(1/j<\eta/20\).
Then fix \(k\) so that both terms in the maximum defining
\(\rho_{j,k}\) are below \(\eta/20\).
This is possible since \(t\ge k\) and \(g(k)/k\to0\).
Thus \(\rho_{j,k}<\eta/10\); the fifth-power bound is
\(2^{\eta N}\) for all sufficiently large \(N\), after absorbing its
fixed constants and polynomial factors. The same pair works at all
these lengths. The machine need not recognize a good pair or know a
convergence modulus for \(g(k)/k\).
For precision, normalize first to a canonical formula on its occurring
variables. If none occur, decide the resulting constant formula
directly, without invoking sparsification or the racing machine.
Otherwise, if their number is \(1\le N'\le N\), its length is polynomial
in \(N'\). For every desired rate the same fixed pair bounds all
sufficiently large \(N'\); the finitely many canonical formulas with
smaller \(N'\) have a common finite maximum running time. Thus the
normalized solver has a uniform \(2^{o(N)}\) bound even when the original
encoding names many unused variables. Parsing and normalization cost
one fixed polynomial in the original input length, so the complete
algorithm contradicts the assumed uniform subexponential lower bound.
\end{proof}

The same model qualification applies: simulation replaces \(g(k)\) by
at most \(\max\{qg(k),2q\}+1=o(k)\) when \(q\) is fixed.
For this uniform transfer, the simulation must itself be uniform in \(k\).
An ineffective family of separately supplied fixed-\(k\) solvers would
not define \(U\); Proposition~\ref{prop:uniform} makes no assertion
about such a family under its weaker premise.

Given two \(n\)-vertex graphs, consider deciding whether their
\(k\)-tuple color histograms agree after two joint rounds, or after
one separate round, starting from atomic types. This includes separate
runs with common color names and pure-tuple comparison in the
disjoint-union run. On the graph pairs produced by the reduction,
Theorem~\ref{thm:choice-reduction} gives equality at every round when
there is no successful choice, and inequality by the specified round
when a choice exists. The inequality persists because each update
retains the old color. Thus the reductions proving
Theorem~\ref{thm:eth} and Proposition~\ref{prop:uniform} give the same
deterministic time exclusions for these fixed-round decisions, with
their respective hypotheses and parameter ranges. The difficulty is
already present in comparing these early histograms; it does not require
many rounds of stabilization.
